For all 27 score/seed combinations, maximizing minimum row or demand coverage gives the same largest feasible threshold within a fixed ranking. Explicit selection \emph{between} the seven original families instead chooses direct excess for row coverage and demand-normalized error for demand coverage in every seed. On the consumed guardian, the demand objective covers 4.94 percentage points more demand and 13.74 points fewer rows (Table~\ref{tab:objectives}). The first-seed descriptive demand-coverage contrast has interval $[2.72,7.02]$ points; it is not a new confirmatory test. This distinguishes an objective comparison from merely reporting demand as a side metric.

The four-way capacity ablation also separates the two heads. Increasing only the error head from 110 to 220 trees lowers guardian row coverage by 0.44--0.63 points, depending on the demand head. Increasing only the demand head raises coverage by 0.20--0.38 points. Joint excess-label MSE rises from 1.0630 to 1.0868 for 110/110 versus 220/220 trees; the error--demand cross moment partly offsets the individual-head errors. Thus demand-head overfitting alone does not explain the smaller model's advantage (Appendix Table~\ref{tab:headfactor}).

Squared, Poisson, and Tweedie losses all target a conditional mean under their moment assumptions. Nevertheless, their finite fits differ: Poisson/Tweedie demand-normalized scores cover 89.40/89.51\% of guardian demand versus 88.23\% for squared loss, at accepted WAPE 0.6296/0.6303 versus 0.6219. Floors also matter: changing 0.01 to 1 shifts first-seed demand-normalized row coverage from 62.74\% to 86.57\%. Hobbies truth-only substitution opens 76.38\% row coverage when replacing the error head, but none when replacing demand alone. These diagnostics identify score and information sensitivity, not a learned solution to group failure; full controls appear in Appendix~\ref{app:objective}.
